\documentclass[aspectratio=169]{beamer}

\usetheme{colorful-dream}

\usepackage{minted}
\setminted{fontsize=\footnotesize, breaklines=true}

\title{Podstawy Pythona dla Data Science}
\subtitle{Wykład 3: Pandas, SciPy i scikit-learn w analizie danych inżynierskich}
\author{Lu Pa}
\date{\today}

% --- Title Slide Logos ---
\titlegraphic{
  \vspace{0.5cm}
  \begin{center}
    \includegraphics[height=1.8cm]{logo_university_placeholder.png}
    \hspace{1cm}
    \includegraphics[height=1.8cm]{logo_faculty_placeholder.png}
  \end{center}
}

\begin{document}

% ---------------------------------------------------------
\begin{frame}
  \titlepage
\end{frame}

% ---------------------------------------------------------
\section{PART I}
\begin{frame}[plain]
  \begin{center}
    \vspace{2cm}
    {\usebeamerfont{title}\color{white}\Huge PART I}
  \end{center}
\end{frame}

% ---------------------------------------------------------
\begin{frame}{Czym jest Pandas?}
  \begin{itemize}
    \item Wysokopoziomowa biblioteka do pracy z danymi tabelarycznymi.
    \item Zbudowana na NumPy.
    \item Zapewnia:
      \begin{itemize}
        \item struktury DataFrame i Series
        \item szybkie filtrowanie i grupowanie
        \item obsługę braków danych
        \item łączenie i scalanie danych
      \end{itemize}
    \item Standard w czyszczeniu i raportowaniu danych.
  \end{itemize}
\end{frame}

% ---------------------------------------------------------
\begin{frame}[fragile]{Wczytywanie danych}
\begin{minted}{python}
import pandas as pd

df = pd.read_csv("../datasets/vehicle_speeds_trackA.csv")
df.info()
df.head()
\end{minted}
\end{frame}

% ---------------------------------------------------------
\begin{frame}{Czyszczenie danych}
  \begin{itemize}
    \item Usuwanie wartości niemożliwych:
      \begin{itemize}
        \item ujemna prędkość
        \item skoki RPM
      \end{itemize}
    \item Obsługa braków danych:
      \begin{itemize}
        \item usuwanie wierszy
        \item interpolacja
      \end{itemize}
    \item Konwersja jednostek (km/h → m/s).
  \end{itemize}
\end{frame}

% ---------------------------------------------------------
\begin{frame}[fragile]{Filtrowanie za pomocą masek logicznych}
\begin{minted}{python}
high_speed = df[df["speed_kmh"] > 80]
braking = df[df["accel_mps2"] < -1.5]
cruise = df[(df["speed_kmh"] > 40) & (df["accel_mps2"].abs() < 0.2)]
\end{minted}
\end{frame}

% ---------------------------------------------------------
\begin{frame}[fragile]{Grupowanie i agregacja}
\begin{minted}{python}
df["minute"] = (df["time_s"] // 60).astype(int)

summary = df.groupby("minute").agg({
    "speed_kmh": ["mean", "max"],
    "fuel_rate_gps": "mean",
    "engine_rpm": "mean"
})
\end{minted}
\end{frame}

% ---------------------------------------------------------
\begin{frame}{SciPy w sygnałach inżynierskich}
  \begin{itemize}
    \item SciPy rozszerza NumPy o:
      \begin{itemize}
        \item przetwarzanie sygnałów
        \item optymalizację
        \item interpolację
        \item statystykę
      \end{itemize}
    \item Dziś: filtrowanie niskoprzepustowe.
  \end{itemize}
\end{frame}

% ---------------------------------------------------------
\begin{frame}[fragile]{Filtr dolnoprzepustowy SciPy}
\begin{minted}{python}
from scipy.signal import butter, filtfilt

b, a = butter(3, 0.05)
v_filtered = filtfilt(b, a, df["speed_kmh"])
\end{minted}
\end{frame}

% ---------------------------------------------------------
\begin{frame}{scikit-learn w modelowaniu}
  \begin{itemize}
    \item Biblioteka ML do:
      \begin{itemize}
        \item regresji
        \item klasyfikacji
        \item klasteryzacji
      \end{itemize}
    \item Dziś: regresja liniowa na danych telemetrycznych.
  \end{itemize}
\end{frame}

% ---------------------------------------------------------
\begin{frame}[fragile]{Prognozowanie zużycia paliwa na podstawie prędkości}
\begin{minted}{python}
from sklearn.linear_model import LinearRegression
import numpy as np

X = df[["speed_kmh"]].values
y = df["fuel_rate_gps"].values

model = LinearRegression()
model.fit(X, y)

print(model.coef_, model.intercept_)
\end{minted}
\end{frame}

% ---------------------------------------------------------
\begin{frame}[fragile]{Wizualizacja dopasowania regresji}
\begin{minted}{python}
plt.scatter(df["speed_kmh"], df["fuel_rate_gps"], s=5)
plt.plot(df["speed_kmh"], model.predict(X), color="red")
plt.xlabel("Prędkość [km/h]")
plt.ylabel("Zużycie paliwa [g/s]")
plt.grid(True)
plt.show()
\end{minted}
\end{frame}

% ---------------------------------------------------------
\section{PART II}
\begin{frame}[plain]
  \begin{center}
    \vspace{2cm}
    {\usebeamerfont{title}\color{white}\Huge PART II}
  \end{center}
\end{frame}

% ---------------------------------------------------------
\begin{frame}{Zadania inżynierskie}
  \begin{itemize}
    \item Wyczyść dane Track A i Track B.
    \item Zidentyfikuj fazy: przyspieszanie, hamowanie, jazda ustalona.
    \item Zastosuj filtr dolnoprzepustowy do prędkości.
    \item Zbuduj model regresji przewidujący zużycie paliwa.
    \item Porównaj styl jazdy na Track A i Track B.
  \end{itemize}
\end{frame}

% ---------------------------------------------------------
\end{document}
